% Bloque C: Isaac, 3:00. S8--S10. Modelos generativos: dos modos, control de
% fidelidad G8 del modo intermedio, SDXL con LoRA y datos BG60k.
\section{Modelos generativos}
\begin{frame}{Dos modos, dos formas de conservar el producto}
  \centering\small
  \begin{tabular}{@{}p{2.6cm}p{5.2cm}p{5.2cm}@{}}
    \toprule
    & {\color{colVenta}\textbf{Intermedio}} & {\color{colCompra}\textbf{Avanzado}} \\
    \midrule
    Modelo & Qwen-Image-2.1-Turbo, edita con instrucciones & SDXL inpainting de 9 canales + LoRA propio \\
    Qué cambia & Toda la imagen, incluida la luz del producto & Solo el fondo; el producto se vuelve a pegar \\
    Ventaja & Mejor integración visual & Píxeles del producto protegidos \\
    Riesgo & Puede redibujar letras o detalles & La unión con el fondo puede verse artificial \\
    GPU (Modal) & H100 / A100 de 80\,GB & L4 / A10G \\
    \bottomrule
  \end{tabular}
  \takeaway{Integración visual o fidelidad: cada modo elige un lado distinto.}
  \note{Isaac. 6:00--7:00 (60 s). Tenemos dos modos. El intermedio usa
  Qwen-Image, un modelo que edita una imagen a partir de un texto. Cambia
  toda la imagen, así que adapta la luz del producto a la escena y se ve
  más integrado, pero puede redibujar una letra o un detalle. El avanzado
  usa SDXL con inpainting: solo rellena el fondo y después vuelve a pegar
  los píxeles originales del producto. El producto queda intacto, pero la
  unión con el fondo puede verse menos natural. Ojo: como cambian a la vez
  el modelo y la forma de conservar el producto, no podemos atribuir una
  diferencia entre modos solo al LoRA.}
\end{frame}

\begin{frame}{Modo intermedio: generar varias candidatas y verificar (G8)}
  \begin{columns}[c]
    \begin{column}{0.5\textwidth}\small
      \textbf{Por cada estilo, hasta 3 candidatas.}\par\medskip
      Se alinea el producto con puntos \textbf{ORB} y \textbf{RANSAC}, y se miden
      las correspondencias, los bordes y el cambio de fondo.\par\medskip
      Se elige la válida que maximiza $r_{\mathrm{inliers}}\cdot\max(0,\rho_{\mathrm{bordes}})$.\par\medskip
      Si falta un estilo, se reintenta dentro de un presupuesto de
      \textbf{420\,s}; si aún falta, falla el lote completo.
    \end{column}
    \begin{column}{0.46\textwidth}
      \begin{block}{Una candidata es válida si}
        \small
        está alineada\\
        $r_{\mathrm{inliers}} \geq 0.45$ \quad (correspondencias)\\
        $\rho_{\mathrm{bordes}} \geq 0.40$ \quad (bordes de Sobel)\\
        $\Delta_{\mathrm{fondo}} \geq 25$ \quad (de 0 a 255)
      \end{block}
    \end{column}
  \end{columns}
  \takeaway{Siempre cinco imágenes verificadas, o un error claro.}
  \source{Pasar los umbrales no asegura que el texto, la marca o las propiedades físicas sean exactos.}
  \note{Isaac. 7:00--8:00 (60 s). Como Qwen puede cambiar el producto, no
  confiamos en una sola salida. Por cada estilo generamos hasta tres
  candidatas y las verificamos: alineamos el producto de la foto original
  con el de la candidata usando puntos ORB y RANSAC, y medimos cuántas
  correspondencias son consistentes, cuánto se parecen los bordes y si el
  fondo realmente cambió. Entre las válidas elegimos la mejor. Si un
  estilo no tiene ninguna válida, reintentamos mientras quepa en 420
  segundos; si igual falta, el lote falla entero, porque prometemos cinco.
  Un límite: estas métricas no garantizan que una etiqueta o un texto
  queden exactos.}
\end{frame}

\begin{frame}{Modo avanzado: SDXL con un LoRA propio entrenado en BG60k}
  \begin{columns}[T]
    \begin{column}{0.48\textwidth}\small
      \textbf{\color{colAcento}LoRA}: adapta el modelo con una actualización de bajo rango
      \[ W' = W_0 + sBA,\qquad r \ll \min(k,d) \]
      Los pesos base quedan congelados.\par\medskip
      \textbf{\color{colAcento}Datos}: piloto curado de BG60k, sin repetir un producto entre particiones.\par\smallskip
      \centering
      \stat{209}{entrenamiento}\hfill\stat{23}{validación}\hfill\stat{25}{prueba}
    \end{column}
    \begin{column}{0.48\textwidth}\small
      \textbf{\color{colAcento}Selección}: revisión ciega, 15 pares\par\smallskip
      \begin{tabular}{@{}lcc@{}}
        \toprule
        & v4 & \textbf{v5 (elegido)} \\
        \midrule
        Preferido & 2 & \textbf{8} \\
        Sin preferencia & \multicolumn{2}{c}{5} \\
        Con equipo de estudio visible & 13/15 & \textbf{0/15} \\
        \bottomrule
      \end{tabular}\par\medskip
      {\scriptsize Cambian a la vez checkpoint y prompts, y los evaluadores no fueron una muestra independiente: no es una prueba con usuarios.}
    \end{column}
  \end{columns}
  \source{BG60k solo permite uso académico.}
  \note{Isaac. 8:00--9:00 (60 s). El modo avanzado usa SDXL de nueve canales,
  que recibe la imagen, la máscara y el ruido, y le agregamos un LoRA
  propio. LoRA entrena solo una actualización pequeña, de bajo rango,
  sobre los pesos del modelo, que quedan congelados. Lo entrenamos en Colab
  con un piloto curado de BG60k: 257 registros, 209 de entrenamiento, 23
  de validación y 25 de prueba, cuidando que un mismo producto no aparezca
  en dos particiones. Para elegir el adaptador hicimos una revisión ciega
  de 15 pares: la versión 5 ganó 8 veces, la 4 dos, y en cinco no hubo
  preferencia. Además, la 4 dejaba ver equipo de estudio en 13 de 15
  imágenes y la 5 en ninguna. Es una evidencia interna, no una prueba con
  usuarios. Nicolás sigue con el agente.}
\end{frame}
